The Cost-Deterred Hiding and Partial Pooling case is the central result of this section, and the mechanism is the one that Section~\ref{ss:renov_problem} makes visible. In Partial Pooling the marginal informed low-quality seller is exactly indifferent between hiding and disclosing a low rating, so $h=0$ and her best option without renovating is worth a disclosed low rating. Renovating gains her $b$ and nothing else. The hiding cost has cancelled, because the equilibrium hiding rate adjusts until the net gain from hiding is zero. In Cost-Deterred Hiding she does not hide and the same comparison applies. Raising the hiding cost therefore moves the hiding rate along the indifference condition and leaves the threshold where it was: tightening enforcement against concealment cannot buy renovation in these regimes, and renovation requires a subsidy that improves the renovate-and-disclose payoff. Complete Adverse Selection is the textbook case. When every informed low-quality seller hides and hiding pays strictly, renovating means giving up the rent $h$, the threshold falls to $b-h$, and enforcement, by lowering $h$, raises the threshold one for one. The sign of $\partial\bar R/\partial(\rho\omega_1)$ in Complete Adverse Selection is the sign of $1-\tilde e$, which is negative whenever silence sells above a verified efficient rating, as it does in apartments.
